\chapter{The Analysis Toolbox}
\label{ch:analysis}

Every tool in this chapter opens from the \menu{Analysis} menu. Most
compute as soon as they open, run on a worker thread so the interface stays
responsive, and export their data as CSV or whitespace-separated
\file{.dat} --- the format is chosen by the extension you type.

\section{Radial distribution function}
\label{sec:rdf}

\menu{Analysis}\menusep\menu{Radial Distribution Function} computes $g(r)$,
the probability of finding an atom at distance $r$ relative to a uniform
density.

\begin{description}
  \item[\ui{Element pair}] Two combos, each \ui{Any} or a specific element.
    \ui{Any}--\ui{Any} gives the total RDF; a specific pair gives the
    partial.
  \item[\ui{r max}] Default 10\,\AA.
  \item[\ui{Bins}] Default 200.
  \item[\ui{Periodic boundary conditions}] Enabled and checked when the
    structure has a cell.
  \item[\ui{Frames}] For trajectories: \ui{start}, \ui{end} and
    \ui{stride}, giving a frame-averaged $g(r)$ over any sub-range.
\end{description}

\begin{calnote}
With periodicity enabled, all periodic images within $r_\text{max}$ are
enumerated explicitly rather than relying on the minimum-image convention.
The result is therefore valid beyond $L/2$ and correct for triclinic cells
--- both places where a naive implementation quietly produces artefacts.
\end{calnote}

\section{Bond length and angle distributions}

\menu{Analysis}\menusep\menu{Bond Length / Angle Distributions} histograms
either pair distances or three-body angles $j$--$i$--$k$ within a cutoff.

\begin{description}
  \item[\ui{Distribution}] \ui{Bond length distribution} or
    \ui{Bond angle distribution (j--i--k)}.
  \item[\ui{Species (pair / center--neighbor)}] Two element filters. For
    angles the first is the central atom.
  \item[\ui{Cutoff}] Default 3\,\AA.
  \item[\ui{Bins}] Default 90.
\end{description}

Angles are in degrees, lengths in \AA{}; both handle periodic images
exactly.

\section{Coordination numbers (CN / GCN)}
\label{sec:coordination}

\menu{Analysis}\menusep\menu{Coordination Numbers (CN / GCN)} computes
per-atom coordination and generalised coordination numbers.

\begin{description}
  \item[\ui{Neighbor cutoff}] \ui{Covalent radii $\times$ tolerance}
    (default tolerance 1.15) or \ui{Fixed cutoff radius} (default 3\,\AA).
  \item[\ui{Bulk CN reference}] The $\mathrm{cn_{max}}$ normalising the GCN
    --- 12 for fcc/hcp, 8 for bcc, 4 for diamond, or
    \ui{Auto (max CN found)}.
\end{description}

Results appear as a per-atom table (index, element, CN, GCN) with a summary
line giving the ranges and means. Two buttons push the result onto the
structure: \ui{Color Viewport by CN} and \ui{Color Viewport by GCN}.

\begin{calnote}
The generalised coordination number weights each neighbour by its own
coordination, which distinguishes sites that plain CN cannot: a terrace
atom, a step-edge atom and a vertex atom on a nanoparticle can share the
same CN but have clearly different GCN. This is what makes GCN such a good
descriptor for catalytic activity.

Periodic images are enumerated explicitly, so a primitive fcc cell
correctly reports CN\,=\,12 even though all twelve neighbours are images of
the single atom in the cell.
\end{calnote}

\section{Static structure factor}

\menu{Analysis}\menusep\menu{Structure Factor S(q)} computes $S(q)$ from
the (frame-averaged) pair distribution via a Lorch-windowed Fourier
transform.

\begin{description}
  \item[\ui{q range}] Defaults 0.3 to 12\,\AA$^{-1}$.
  \item[\ui{q points}] Default 400.
  \item[\ui{g(r) r max}] Upper bound of the Fourier integral, default
    10\,\AA{}. Larger values improve low-$q$ fidelity.
  \item[\ui{g(r) bins}] Default 400.
  \item[\ui{Frames}] Start/end/stride, as for the RDF.
\end{description}

\section{X-ray diffraction}

\menu{Analysis}\menusep\menu{X-Ray Diffraction (XRD)} simulates a powder
pattern with the Debye scattering equation.

\begin{description}
  \item[\ui{Wavelength}] Presets for Cu\,K$\alpha$ (1.54056\,\AA),
    Co\,K$\alpha$, Mo\,K$\alpha$, Cr\,K$\alpha$, or \ui{Custom} to type
    your own.
  \item[\ui{2$\theta$ range}] Defaults 10--90\textdegree.
  \item[\ui{Points}] Default 800.
  \item[\ui{Supercell repeat}] Periodic structures are repeated before the
    Debye sum, default 3. More repeats sharpen the Bragg peaks, but the
    cost grows as $N^2$ in atoms.
\end{description}

Unlike the other analysis dialogs, this one waits for you to press
\ui{Simulate}. The status line then reports which form factors were used:
Waasmaier-Kirfel factors when every species is tabulated, otherwise a
constant $f = Z$ approximation --- in which case peak \emph{positions}
remain exact but high-angle \emph{intensities} are approximate.

Two exports: \ui{Export Curve\ldots} writes the $2\theta$--intensity
pattern, and \ui{Export Peaks\ldots} writes detected peaks (those above
2\,\% of the maximum) with their $d$-spacings from
$d = \lambda / (2\sin\theta)$.

\section{Warren-Cowley chemical order}
\label{sec:warren-cowley}

\menu{Analysis}\menusep\menu{Warren-Cowley analysis} quantifies chemical
short-range order in multicomponent alloys through
\[
  \alpha_{ij} = 1 - \frac{p_{ij}}{c_j},
\]
where $p_{ij}$ is the probability that a neighbour of an $i$-type atom is
of type $j$, and $c_j$ is the overall concentration of $j$.

\begin{description}
  \item[$\alpha = 0$] The ideal random alloy.
  \item[$\alpha < 0$] Ordering --- unlike pairs preferred.
  \item[$\alpha > 0$] Clustering or segregation of like species.
\end{description}

Set \ui{First shell cutoff} (default 3.2\,\AA) and, optionally,
\ui{Second shell cutoff} (default 4.8\,\AA). The dialog reports the
concentrations, the mean coordination of each shell, and the full
$\alpha_{ij}$ matrix for every ordered species pair, exportable as CSV.

\begin{caltip}
A quick sanity check: a perfectly ordered B2 (CsCl-type) binary gives
$\alpha_{AB} = -1$ in the first shell and $+1$ in the second, while a
random decoration of the same lattice gives values near zero.
\end{caltip}

\section{Local entropy}

\menu{Analysis}\menusep\menu{Local Entropy Analysis} computes the per-atom
pair-entropy fingerprint of Piaggi and Parrinello, in units of $k_B$:
\[
  s_S^i = -2\pi\rho \int_0^{r_c}
    \bigl[\, g_i(r)\ln g_i(r) - g_i(r) + 1 \,\bigr] r^2\,\mathrm{d}r ,
\]
with $g_i(r)$ the Gaussian-smoothed radial distribution around atom $i$.

\begin{description}
  \item[\ui{Cutoff}] Integration cutoff, default 5\,\AA{} --- three or four
    coordination shells works well.
  \item[\ui{Broadening $\sigma$}] Gaussian width, default 0.15\,\AA.
  \item[\ui{Radial grid points}] Default 100.
  \item[\ui{Average over neighbors}] Smooths $s_i$ over each atom's
    neighbourhood, sharpening the contrast between ordered and disordered
    regions.
\end{description}

The result is stored as the per-atom field \texttt{local\_entropy} and
mapped onto the atoms immediately, with a distribution histogram in the
dialog. \emph{Lower} (more negative) values mark more ordered,
crystal-like environments; higher values mark disordered, liquid-like ones.
This makes it an effective order parameter for melting, solid-liquid
interfaces and grain boundaries.

\section{Raman modes}

\menu{Analysis}\menusep\menu{Raman Modes} performs a $\Gamma$-point
factor-group analysis: which optical phonons are Raman-active, IR-active or
silent, by symmetry alone.

The method uses no hard-coded character tables. spglib supplies the space
group operations of the primitive cell; the character table of the isogonal
point group is computed numerically from the class-sum algebra; the
mechanical representation $\chi(R) = (\pm 1 + 2\cos\theta)\,N_\text{unmoved}(R)$
is reduced into irreducible representations; the acoustic translations are
subtracted; and activity follows from the vector (IR) and symmetric
polarizability (Raman) representations.

The dialog reports the space group, point group, number of atoms in the
primitive cell, and a table of irreps with their degeneracy, optical mode
count and activity.

\begin{caltip}
For silicon in the diamond structure the result is the textbook
$\Gamma = T_{2g} + T_{1u}$: one triply degenerate Raman-active optical
mode, and the acoustic branch.
\end{caltip}

\begin{calnote}
Mulliken subscripts are assigned from Cartesian axis geometry. In
low-symmetry orthorhombic groups the subscript convention can be permuted
relative to a particular textbook's axis choice --- the activities and
degeneracies are unaffected.
\end{calnote}

\begin{calwarn}
Needs \texttt{spglib} and a periodic structure.
\end{calwarn}

\section{Volumetric data}
\label{sec:volumetric}

\menu{Analysis}\menusep\menu{Volumetric Data} visualises 3D scalar fields:
charge densities, electrostatic potentials, wavefunctions, ELF.

\subsection{Loading fields}

\ui{Load Field A\ldots} and \ui{Load Field B\ldots} read Gaussian
\file{.cube} (including files written by Quantum ESPRESSO's \texttt{pp.x}),
VASP \file{CHGCAR} / \file{LOCPOT} / \file{PARCHG} / \file{ELFCAR}, and
XCrySDen \file{.xsf} grids. Each is reported with its grid dimensions and
value range.

Field A defines the geometry; Field B is optional and colours it.

\subsection{Render modes}

\ui{Edit Volumetric Render} offers two modes: \ui{Isosurfaces} and \ui{Color
Slice}. Potential-map colouring used to be a third; it is now a group inside
\ui{Isosurfaces}, because it always was one --- the same surface, extracted the
same way at the same isovalue, differing only in what colours it.

\subsection{Isosurfaces}

The isovalue slider and numeric box re-extract the surface in real time.
Extraction uses a marching-cubes-family algorithm on a tetrahedral
decomposition of each grid cell, which is topologically unambiguous, with
smooth normals from the field gradient and periodic stitching so surfaces
close correctly across cell boundaries.

\subsection{Potential map colour}

The \ui{Potential Map Color} group under \ui{Isosurfaces} makes an EPM: the
isosurface (say the electron density) shapes the geometry, and the field chosen
under \ui{Color by} (say the Hartree potential) is interpolated at every
surface vertex and mapped through the \ui{Color map}. \ui{Invert palette}
flips the ramp.

\ui{Custom range} pins the colour scale to an explicit min/max instead of the
colouring field's own extremes. A potential runs from deeply negative at the
nuclei to nearly flat in the vacuum, so on the full range everything
interesting sits in a sliver of the ramp --- and a fixed window is also what
lets two molecules be compared on one scale.

Because the colouring is an option on the surface rather than a mode, it
applies to \emph{every} ticked dataset at once: two orbitals can be shown side
by side, both painted by the same potential.

\begin{caltip}
The classic recipe: load the charge density as Field A, the local potential
as Field B, pick an isovalue on the density that traces the molecular
surface, and choose \ui{Coolwarm} --- the diverging palette puts the
neutral potential at white, so electron-rich and electron-poor regions read
immediately.
\end{caltip}

\subsection{Slice planes}

\ui{Color Slice} cuts through the volume with a colour-mapped plane, oriented
by Miller indices $(h\,k\,l)$ against the grid's own lattice and swept along its
normal by the \ui{Offset} slider.

\ui{Extent} sets how far the plane is drawn, in unit cells across. \ui{This
unit cell only} keeps it inside the box the field was computed in --- the
honest extent, and the one to use when reading values off it. The replicated
options tile the periodic field over the neighbouring cells, which is what
makes a surface reconstruction or an adsorbate pattern legible instead of
showing one cell floating on its own. \ui{Outline the slice} draws its
boundary, useful where the field has gone to zero and the quad would otherwise
fade into the background.

\ui{Custom color range} pins the ramp the same way the potential map's does:
a few outlier voxels (a nuclear cusp, a spike at a boundary) otherwise compress
everything else into one end of it.

\subsection{Export}

\ui{Export Isosurface (OBJ)\ldots} writes the triangle mesh with normals
for external rendering; \ui{Export Slice (CSV)\ldots} writes the sampled
plane as $x, y, z$, value.

\section{Charge density difference}
\label{sec:cdd}

\menu{Analysis}\menusep\menu{Charge Density Difference (CDD)\dots} computes
\[
  \Delta\rho \;=\; \rho_{A+B} \;-\; \rho_{A} \;-\; \rho_{B},
\]
the redistribution of charge caused by putting two fragments together. Each
of the three densities on its own is dominated by the atomic cores and shows
nothing; the difference shows the bond.

\subsection{Step 1 --- density source}

Pick a completed \menu{Simulation}\menusep\menu{Single-point Calculation} that
saved its wavefunctions (GPAW \file{.gpw}); that run supplies $\rho_{A+B}$ and,
more importantly, the calculator both fragments are rebuilt from. Choose
whether to difference the \ui{All-electron density} or the
\ui{Pseudodensity}. The pseudodensity is usually the clearer field to read:
it carries no nuclear cusps, so the isosurface scale is set by the bonding
features rather than by the cores.

\subsection{Step 2 --- subsystem partition}

Two columns list the atoms of the selected run. Move atoms across with the
arrow buttons or by double-clicking a row. The split is exhaustive --- every
atom belongs to exactly one side --- because $\Delta\rho$ only integrates to
zero if $A$ and $B$ together are the whole system.

\subsection{What the run does}

Both fragments are recomputed in the parent's own unit cell, at the parent's
geometry, with the parameters read back out of its \file{.gpw}: cutoff, XC
functional, grid, $k$-points and convergence cannot drift between the three
terms. \emph{Nothing is relaxed.} $\Delta\rho$ is defined at one geometry, and
letting a fragment relax would mix charge transfer with structural
rearrangement into a quantity nobody can read off an isosurface.

\begin{calnote}
Cutting a closed-shell molecule in half leaves open-shell fragments, and an
isolated atom with a partly-filled degenerate shell often will not converge
under the settings that converged the molecule in one pass. The run therefore
tries the parent's exact parameters first, then adds occupation smearing, then
spin polarization, and reports in the log which was needed. The reported
\texttt{net} charge is the check: $A$ and $B$ together hold exactly the
electrons $A{+}B$ does, so it must come out at zero.
\end{calnote}

\subsection{Result}

The difference is written as \file{cdd.cube} and loaded automatically into the
\ui{Volumetric Data} dock (\S\ref{sec:volumetric}), where it renders in the
main 3D viewport like any other grid. A \ui{Coolwarm} colormap is the natural
choice: accumulation and depletion sit on opposite sides of white. The status
bar reports how much charge was redistributed, from \file{cdd.json}.

\section{Adsorption and catalysis}
\label{sec:adsorption}

\menu{Analysis}\menusep\menu{Adsorption \& Catalysis} finds high-symmetry
adsorption sites on a slab and populates them.

\subsection{Site detection}

Opening the dialog detects sites on the top surface automatically and
reports a census, for example \emph{4 top, 12 bridge, 4 fcc, 4 hcp}:

\begin{description}
  \item[\texttt{top}] Directly above a surface atom.
  \item[\texttt{bridge}] Midway between nearest-neighbour surface atoms.
  \item[\texttt{fcc}, \texttt{hcp}] Threefold hollows, distinguished by
    whether a second-layer atom sits directly underneath (hcp) or not
    (fcc).
  \item[\texttt{hollow}] A threefold hollow that could not be classified,
    typically because there is no second layer.
\end{description}

Neighbour vectors are enumerated over periodic images, so the census is
correct even for a small $p(1\times1)$ cell where surface atoms bond to
their own periodic copies. The table lists every site with its in-plane
coordinates and can be filtered by type.

\subsection{Placing adsorbates}

\begin{description}
  \item[\ui{Adsorbate}] \texttt{OH}, \texttt{O}, \texttt{H}, \texttt{CO},
    \texttt{CHO}, \texttt{H2O}, \texttt{N2}, \texttt{O2}, \texttt{NH3},
    \texttt{CH4}, or any \texttt{ase.build.molecule} name or chemical
    formula you type.
  \item[\ui{Height}] Anchor-atom height above the surface layer, default
    2\,\AA.
  \item[\ui{Place on Selected Sites}] Adds one copy per selected table row.
\end{description}

Molecules are anchored by the chemically sensible atom and oriented
upright: O down for \texttt{OH} and \texttt{H2O}, C down for \texttt{CO}
and \texttt{CHO}, N down for \texttt{NH3}.

\subsection{Coverage series}

The \ui{Coverage series} group generates a systematic set in one step:
choose a \ui{Site family} and check the coverages you want --- 0.25, 0.50,
0.75, 1.00\,ML. Calango places evenly strided subsets of that site family
and opens one labelled tab per coverage, ready for a batch of adsorption
energy calculations.

\section{Brillouin zone and k-path builder}
\label{sec:bz}

\menu{Analysis}\menusep\menu{Brillouin Zone / k-Path} shows the first
Brillouin zone as a Wigner-Seitz cell of the reciprocal lattice, with
high-symmetry points labelled from ASE's Bravais lattice detection.

Drag to rotate, \kbd{Shift}+drag to pan, wheel to zoom;
\ui{Orthographic projection} removes perspective foreshortening when you
need to read the zone geometry.

\subsection{Building a path}

Click high-symmetry points in the 3D view to append them to the path. The
list shows the sequence with fractional coordinates.

\begin{description}
  \item[\ui{Suggested}] Loads ASE's suggested path for the lattice.
  \item[\ui{Break}] Starts a new discontinuous section, giving paths like
    $\Gamma \rightarrow X \,|\, M \rightarrow R$.
  \item[\ui{Undo}, \ui{Clear}] Remove the last point, or start over.
  \item[\ui{Points per segment}] Sampling density, default 40.
\end{description}

\subsection{Export}

\ui{Export k-Path\ldots} writes the path for an external code:

\begin{center}
\small
\begin{tabular}{@{}ll@{}}
\toprule
Format & Typical file \\
\midrule
VASP \texttt{KPOINTS} (line mode)          & \file{KPOINTS} \\
Quantum ESPRESSO \texttt{K\_POINTS crystal\_b} & \file{kpath\_qe.in} \\
CASTEP \texttt{SPECTRAL\_KPOINT\_PATH}     & \file{kpath\_castep.cell} \\
SIESTA \texttt{BandLines}                  & \file{kpath\_siesta.fdf} \\
ASE / Python script                        & \file{kpath\_ase.py} \\
\bottomrule
\end{tabular}
\end{center}

\ui{Export Figure (PNG/SVG)\ldots} renders the zone at high resolution;
the SVG output is vector and drops straight into a paper figure.
